\documentclass[10pt,a4paper]{amsart}
\usepackage{amsmath,amssymb,amsthm,graphicx}
% Retained minimal preamble for the Stacks Project chapter varieties.tex.
% Verbatim excerpt of the chapter preamble: Xy-pic setup (native preamble lines 14-18), the native theorem
% declarations with the mathematical macros (native preamble lines 154-217), and the project own no-op
% definitions of the reference/slogan/history display environments (native wrapper.tex lines 1-3) that the
% counted statement uses for its slogan. Copied mechanically; nothing added, removed or redefined.
% For commutative diagrams we use Xy-pic
\usepackage[all]{xy}

% We use 2cell for 2-commutative diagrams.
\xyoption{2cell}
% Theorem environments.
%
\theoremstyle{plain}
\newtheorem{theorem}[subsection]{Theorem}
\newtheorem{proposition}[subsection]{Proposition}
\newtheorem{lemma}[subsection]{Lemma}

\theoremstyle{definition}
\newtheorem{definition}[subsection]{Definition}
\newtheorem{example}[subsection]{Example}
\newtheorem{exercise}[subsection]{Exercise}
\newtheorem{situation}[subsection]{Situation}

\theoremstyle{remark}
\newtheorem{remark}[subsection]{Remark}
\newtheorem{remarks}[subsection]{Remarks}

\numberwithin{equation}{subsection}

% Macros
%
\def\lim{\mathop{\mathrm{lim}}\nolimits}
\def\colim{\mathop{\mathrm{colim}}\nolimits}
\def\Spec{\mathop{\mathrm{Spec}}}
\def\Hom{\mathop{\mathrm{Hom}}\nolimits}
\def\Ext{\mathop{\mathrm{Ext}}\nolimits}
\def\SheafHom{\mathop{\mathcal{H}\!\mathit{om}}\nolimits}
\def\SheafExt{\mathop{\mathcal{E}\!\mathit{xt}}\nolimits}
\def\Sch{\mathit{Sch}}
\def\Mor{\mathop{\mathrm{Mor}}\nolimits}
\def\Ob{\mathop{\mathrm{Ob}}\nolimits}
\def\Sh{\mathop{\mathit{Sh}}\nolimits}
\def\NL{\mathop{N\!L}\nolimits}
\def\CH{\mathop{\mathrm{CH}}\nolimits}
\def\proetale{{pro\text{-}\acute{e}tale}}
\def\etale{{\acute{e}tale}}
\def\QCoh{\mathit{QCoh}}
\def\Ker{\mathop{\mathrm{Ker}}}
\def\Im{\mathop{\mathrm{Im}}}
\def\Coker{\mathop{\mathrm{Coker}}}
\def\Coim{\mathop{\mathrm{Coim}}}

% Boxtimes
%
\DeclareMathSymbol{\boxtimes}{\mathbin}{AMSa}{"02}

%
% Macros for moduli stacks/spaces
%
\def\QCohstack{\mathcal{QC}\!\mathit{oh}}
\def\Cohstack{\mathcal{C}\!\mathit{oh}}
\def\Spacesstack{\mathcal{S}\!\mathit{paces}}
\def\Quotfunctor{\mathrm{Quot}}
\def\Hilbfunctor{\mathrm{Hilb}}
\def\Curvesstack{\mathcal{C}\!\mathit{urves}}
\def\Polarizedstack{\mathcal{P}\!\mathit{olarized}}
\def\Complexesstack{\mathcal{C}\!\mathit{omplexes}}
% \Pic is the operator that assigns to X its picard group, usage \Pic(X)
% \Picardstack_{X/B} denotes the Picard stack of X over B
% \Picardfunctor_{X/B} denotes the Picard functor of X over B
\def\Pic{\mathop{\mathrm{Pic}}\nolimits}
\def\Picardstack{\mathcal{P}\!\mathit{ic}}
\def\Picardfunctor{\mathrm{Pic}}
\def\Deformationcategory{\mathcal{D}\!\mathit{ef}}
\newenvironment{reference}{}{}
\newenvironment{slogan}{}{}
\newenvironment{history}{}{}
\usepackage{cleveref}
\crefname{originalreference}{original reference}{original references}
\Crefname{originalreference}{Original reference}{Original references}
\makeatletter
\newlabel{morphisms-lemma-closed-immersion-proper}{{01W5}{source}{}{sourcepointer-1}{}}
\newlabel{morphisms-lemma-closed-immersion-proper@cref}{{[originalreference][0][]01W5}{[1][0][]source}{}{}{}}
\makeatother
\crefname{originalreference}{original reference}{original references}
\Crefname{originalreference}{Original reference}{Original references}
\makeatletter
\newlabel{morphisms-lemma-composition-proper}{{01W3}{source}{}{sourcepointer-2}{}}
\newlabel{morphisms-lemma-composition-proper@cref}{{[originalreference][0][]01W3}{[1][0][]source}{}{}{}}
\makeatother
\crefname{originalreference}{original reference}{original references}
\Crefname{originalreference}{Original reference}{Original references}
\makeatletter
\newlabel{morphisms-lemma-quasi-compact-immersion}{{01RG}{source}{}{sourcepointer-3}{}}
\newlabel{morphisms-lemma-quasi-compact-immersion@cref}{{[originalreference][0][]01RG}{[1][0][]source}{}{}{}}
\makeatother
\crefname{originallemma}{original lemma}{original lemmas}
\Crefname{originallemma}{Original lemma}{Original lemmas}
\makeatletter
\newlabel{lemma-degree-birational-pullback}{{0AYU}{source}{}{sourcepointer-4}{}}
\newlabel{lemma-degree-birational-pullback@cref}{{[originallemma][0][]0AYU}{[1][0][]source}{}{}{}}
\makeatother
\crefname{originalreference}{original reference}{original references}
\Crefname{originalreference}{Original reference}{Original references}
\makeatletter
\newlabel{coherent-lemma-power-ideal-kills-sheaf}{{01Y9}{source}{}{sourcepointer-5}{}}
\newlabel{coherent-lemma-power-ideal-kills-sheaf@cref}{{[originalreference][0][]01Y9}{[1][0][]source}{}{}{}}
\makeatother
\crefname{originalreference}{original reference}{original references}
\Crefname{originalreference}{Original reference}{Original references}
\makeatletter
\newlabel{coherent-lemma-i-star-equivalence}{{087T}{source}{}{sourcepointer-6}{}}
\newlabel{coherent-lemma-i-star-equivalence@cref}{{[originalreference][0][]087T}{[1][0][]source}{}{}{}}
\makeatother
\crefname{originalreference}{original reference}{original references}
\Crefname{originalreference}{Original reference}{Original references}
\makeatletter
\newlabel{algebra-lemma-length-additive}{{00IV}{source}{}{sourcepointer-7}{}}
\newlabel{algebra-lemma-length-additive@cref}{{[originalreference][0][]00IV}{[1][0][]source}{}{}{}}
\makeatother
\crefname{originalreference}{original reference}{original references}
\Crefname{originalreference}{Original reference}{Original references}
\makeatletter
\newlabel{algebra-lemma-dimension-is-length}{{00IY}{source}{}{sourcepointer-8}{}}
\newlabel{algebra-lemma-dimension-is-length@cref}{{[originalreference][0][]00IY}{[1][0][]source}{}{}{}}
\makeatother
\crefname{originalreference}{original reference}{original references}
\Crefname{originalreference}{Original reference}{Original references}
\makeatletter
\newlabel{algebra-lemma-length-independent}{{00IX}{source}{}{sourcepointer-9}{}}
\newlabel{algebra-lemma-length-independent@cref}{{[originalreference][0][]00IX}{[1][0][]source}{}{}{}}
\makeatother
\crefname{originallemma}{original lemma}{original lemmas}
\Crefname{originallemma}{Original lemma}{Original lemmas}
\makeatletter
\newlabel{lemma-degree-on-proper-curve}{{0AYV}{source}{}{sourcepointer-10}{}}
\newlabel{lemma-degree-on-proper-curve@cref}{{[originallemma][0][]0AYV}{[1][0][]source}{}{}{}}
\makeatother
\title{The degree of a locally free module on a proper scheme of dimension at most one is the multiplicity-weighted sum of its degrees on the one-dimensional components}
\author{The Stacks Project Authors}
\date{stacks-project:varieties.tex:local-snapshot-20261011; GNU FDL 1.2 or later}
\begin{document}
\maketitle
\noindent\textit{Editorial note.} All delivered mathematical text is the author's own, copied verbatim from the pinned Stacks Project chapter file varieties.tex; no statement or proof step was rewritten, reordered or removed.\par\medskip
\noindent\textit{Editorial note.} The counted span cites ten cross-chapter labels, each recorded in assembly.external\_references with its canonical tag as a navigation pointer.\par\medskip
\noindent\textit{Editorial note.} The proof cites no bibliography item inside the counted span.\par\medskip
\noindent\textit{Allowed standard prerequisites (editorial).}\begin{itemize}
\item Birational invariance of the degree for morphisms inducing isomorphisms at generic points of components. Stacks Project, Varieties, Lemma lemma-degree-birational-pullback (tag 0AYU). Reduces the general statement to the irreducible case via the coproduct morphism.
\item Nilpotency of the ideal of a closed subscheme and length additivity at the generic point. Stacks Project, Cohomology of Schemes, Lemma coherent-lemma-power-ideal-kills-sheaf (tag 01YI); Algebra, Lemmas algebra-lemma-length-additive (tag 00IV) and algebra-lemma-dimension-is-length (tag 00IZ). Supply the nilpotency index and the length computation that drives the induction.
\end{itemize}\medskip
\ifdefined\hypertarget\hypertarget{sourcepointer-1}{}\fi
\noindent\textit{Source reference.} Original reference 01W5: sources/stacks-chow/tags.txt, source lines 2470--2470. Cross-chapter standard prerequisite cited by the author inside the counted span; it is used as a pointer and its own proof is not part of this item.\par\medskip
\ifdefined\hypertarget\hypertarget{sourcepointer-2}{}\fi
\noindent\textit{Source reference.} Original reference 01W3: sources/stacks-chow/tags.txt, source lines 2468--2468. Cross-chapter standard prerequisite cited by the author inside the counted span; it is used as a pointer and its own proof is not part of this item.\par\medskip
\ifdefined\hypertarget\hypertarget{sourcepointer-3}{}\fi
\noindent\textit{Source reference.} Original reference 01RG: sources/stacks-chow/tags.txt, source lines 2301--2301. Cross-chapter standard prerequisite cited by the author inside the counted span; it is used as a pointer and its own proof is not part of this item.\par\medskip
\ifdefined\hypertarget\hypertarget{sourcepointer-4}{}\fi
\noindent\textit{Source reference.} Original lemma 0AYU: sources/stacks-chow/tags.txt, source lines 13802--13802. Cross-chapter standard prerequisite cited by the author inside the counted span; it is used as a pointer and its own proof is not part of this item.\par\medskip
\ifdefined\hypertarget\hypertarget{sourcepointer-5}{}\fi
\noindent\textit{Source reference.} Original reference 01Y9: sources/stacks-chow/tags.txt, source lines 2548--2548. Cross-chapter standard prerequisite cited by the author inside the counted span; it is used as a pointer and its own proof is not part of this item.\par\medskip
\ifdefined\hypertarget\hypertarget{sourcepointer-6}{}\fi
\noindent\textit{Source reference.} Original reference 087T: sources/stacks-chow/tags.txt, source lines 10438--10438. Cross-chapter standard prerequisite cited by the author inside the counted span; it is used as a pointer and its own proof is not part of this item.\par\medskip
\ifdefined\hypertarget\hypertarget{sourcepointer-7}{}\fi
\noindent\textit{Source reference.} Original reference 00IV: sources/stacks-chow/tags.txt, source lines 693--693. Cross-chapter standard prerequisite cited by the author inside the counted span; it is used as a pointer and its own proof is not part of this item.\par\medskip
\ifdefined\hypertarget\hypertarget{sourcepointer-8}{}\fi
\noindent\textit{Source reference.} Original reference 00IY: sources/stacks-chow/tags.txt, source lines 696--696. Cross-chapter standard prerequisite cited by the author inside the counted span; it is used as a pointer and its own proof is not part of this item.\par\medskip
\ifdefined\hypertarget\hypertarget{sourcepointer-9}{}\fi
\noindent\textit{Source reference.} Original reference 00IX: sources/stacks-chow/tags.txt, source lines 695--695. Cross-chapter standard prerequisite cited by the author inside the counted span; it is used as a pointer and its own proof is not part of this item.\par\medskip
\ifdefined\hypertarget\hypertarget{sourcepointer-10}{}\fi
\noindent\textit{Source reference.} Original lemma 0AYV: sources/stacks-chow/tags.txt, source lines 13803--13803. Cross-chapter standard prerequisite cited by the author inside the counted span; it is used as a pointer and its own proof is not part of this item.\par\medskip
\begin{lemma}
\label{lemma-degree-in-terms-of-components}
Let $k$ be a field. Let $X$ be a proper scheme of dimension $\leq 1$ over $k$.
Let $\mathcal{E}$ be a locally free $\mathcal{O}_X$-module of rank $n$.
Then
$$
\deg(\mathcal{E}) = \sum m_i \deg(\mathcal{E}|_{C_i})
$$
where $C_i \subset X$, $i = 1, \ldots, t$ are the irreducible components
of dimension $1$ with reduced induced scheme structure and $m_i$ is the
multiplicity of $C_i$ in $X$.
\end{lemma}

\begin{proof}
Observe that the statement makes sense because $C_i \to \Spec(k)$
is proper of dimension $1$; we will use below that any closed subscheme
of $X$ is proper over $k$ by Morphisms, Lemmas
\ref{morphisms-lemma-closed-immersion-proper} and
\ref{morphisms-lemma-composition-proper}. Consider the open subscheme
$U_i = X \setminus (\bigcup_{j \not = i} C_j)$ and let $X_i \subset X$
be the scheme theoretic closure of $U_i$. Note that $X_i \cap U_i = U_i$
(scheme theoretically) and that $X_i \cap U_j = \emptyset$
(set theoretically) for $i \not = j$; this follows from the description
of scheme theoretic closure in
Morphisms, Lemma \ref{morphisms-lemma-quasi-compact-immersion}.
Thus we may apply Lemma \ref{lemma-degree-birational-pullback} to the morphism
$X' = \coprod X_i \to X$. Since it is clear that $C_i \subset X_i$
(scheme theoretically) and that the multiplicity of $C_i$ in $X_i$ is equal
to the multiplicity of $C_i$ in $X$, we see that we reduce to the case
discussed in the following paragraph.

\medskip\noindent
Assume $X$ is irreducible with generic point $\xi$. Let
$C = X_{red}$ have multiplicity $m$.
We have to show that $\deg(\mathcal{E}) = m \deg(\mathcal{E}|_C)$.
Let $\mathcal{I} \subset \mathcal{O}_X$ be the ideal defining the closed
subscheme $C$. Let $e \geq 0$ be minimal such that $\mathcal{I}^{e + 1} = 0$
(Cohomology of Schemes, Lemma \ref{coherent-lemma-power-ideal-kills-sheaf}).
We argue by induction on $e$. If $e = 0$, then $X = C$ and the result
is immediate. Otherwise we set $\mathcal{F} = \mathcal{I}^e$ viewed as
a coherent $\mathcal{O}_C$-module (Cohomology of Schemes, Lemma
\ref{coherent-lemma-i-star-equivalence}).
Let $X' \subset X$ be the closed subscheme cut out by the
coherent ideal $\mathcal{I}^e$ and let $m'$ be the multiplicity
of $C$ in $X'$. Taking stalks at $\xi$ of the short exact sequence
$$
0 \to \mathcal{F} \to \mathcal{O}_X \to \mathcal{O}_{X'} \to 0
$$
we find (use Algebra, Lemmas \ref{algebra-lemma-length-additive},
\ref{algebra-lemma-dimension-is-length}, and
\ref{algebra-lemma-length-independent}) that
$$
m = \text{length}_{\mathcal{O}_{X, \xi}} \mathcal{O}_{X, \xi}
= \dim_{\kappa(\xi)} \mathcal{F}_\xi +
\text{length}_{\mathcal{O}_{X', \xi}} \mathcal{O}_{X', \xi}
= r + m'
$$
where $r$ is the rank of $\mathcal{F}$ as a coherent sheaf on $C$.
Tensoring with $\mathcal{E}$ we obtain a short exact sequence
$$
0 \to \mathcal{E}|_C \otimes \mathcal{F} \to \mathcal{E} \to
\mathcal{E} \otimes \mathcal{O}_{X'} \to 0
$$
By induction we have
$\deg(\mathcal{E}|_{X'}) = m' \deg(\mathcal{E}|_C)$.
By Lemma \ref{lemma-degree-on-proper-curve} we have
$\chi(\mathcal{E}|_C \otimes \mathcal{F}) =
r \deg(\mathcal{E}|_C) + n \chi(\mathcal{F})$.
Putting everything together we obtain the result.
\end{proof}

\end{document}
